\subsection{The same first failure with the original factorial exponents}
\label{sec:ghost-class3-original}

\begin{result}{C.2.6}{Proved}{An A.44 prime choice first fails integrally at class three}
\label{thm:ghost-original-first-integral-class3}
There is a recursively specified choice of primes permitted
by SECTIONS A.44, with its original exponents
\(n_j=(j+4)!\), for which the ordinary pair
\[
              A=3\prod_jq_j^{n_j},\qquad
              B=-2\prod_jr_j^{n_j}
\]
has first unrestricted integral failure at nilpotency
class three. Every integral class-two section above
this pair has a nonzero local third obstruction at
\(q_0\) in coefficient prime three, already detected
modulo nine.
\end{result}
\begin{proof}
We specify the new finite field and prime test, and check
the coefficient-prime-three issues that are absent from
Theorem~\ref{thm:ghost-first-integral-class3}.

Let \(H_3\) have the same Heisenberg coordinate law
over \(\mathbb F_3\). It has order \(27\), exponent
three, center \(Z_3\) of order three, and arithmetic
action \((a,b,h)\mapsto(\lambda a,\lambda b,\lambda^2h)\)
for \(\lambda\in\mathbb F_3^\times\). The point-three
section defines a surjection onto
\(H_3\rtimes\mathbb F_3^\times\), of order \(54\).
Indeed over \(K_3=\mathbb Q(\mu_3)\), the Kummer
classes of \(3,-2\) are independent: valuations at
three and two are respectively two and one, both
nonzero modulo three. Surjectivity onto the abelian
quotient forces the full Heisenberg image by the
nonzero commutator. Since \(-1-1=1\) in
\(\mathbb F_3\), the commutator subgroup of the full
semidirect product is all of \(H_3\), exactly as in
the coefficient-prime-five calculation. Its field
\(L_3\) therefore has
\[
                         L_3^{\rm ab}=K_3.
\]

An explicit presentation of the \emph{actual} point-three
field is
\[
 \begin{gathered}
 \zeta=\zeta_3,\quad X^3=3,\quad Y^3=-2,\\
 \Theta_3(X)=(1-\zeta X)(1-\zeta^2X)^2,
 \quad Z_0^3=\Theta_3(X),\quad
 L_3=\mathbb Q(\zeta,X,Y,Z_0).
 \end{gathered}
\]
The proof of the generic-cover identification in
Theorem~\ref{thm:ghost-first-integral-class3} applies
with the index range \(0,1,2\), with the following
literal formulas:
\[
 \begin{gathered}
 \Theta_3(\zeta X)/\Theta_3(X)
          =(1-X)^3/(1-X^3),\\
 \sigma(X)=\zeta X,\quad \sigma(Y)=Y,\quad
 \sigma(Z_0)=Z_0(1-X)/Y,\\
 \tau(X)=X,\quad\tau(Y)=\zeta Y,\quad\tau(Z_0)=Z_0.
 \end{gathered}
\]
The two commuting automorphisms of the abelian radical
base could not act this way on a cube root lying in
that base, so adjoining \(Z_0\) has degree three.
The tangential branch has \(Y(0)=Z_0(0)=1\).
The nontrivial cyclotomic descent is
\[
 \zeta\longmapsto\zeta^2,\quad X,Y\longmapsto X,Y,
       \quad Z_0\longmapsto Z_0^2/(1-\zeta^2X).
\]
It preserves that branch and sends \(\sigma,\tau\)
to \(\sigma^2,\tau^2\); its action on their
commutator is the identity. Thus this is precisely
the canonical class-two exponent-three cover with
its tangential arithmetic action, specialized at
three. No choice of a different Heisenberg extension
has entered the prime condition.

In the first A.44 stage take all old local solutions
\(X_p=Y_p=1\). Here \(n_0=24\), \(n_1=120\).
Let \(M\) be its original CRT modulus, the product
of \(p^{H_1(p)}\) over the primes at most
\(B(120)\) and the divisors of six. Define
\[
 D_3=\operatorname{lcm}(M,24,9),\quad
 F_3=\mathbb Q(\mu_{D_3}),\quad
 R_3=F_3\bigl((-2)^{1/24},3^{1/3}\bigr).
\]
The field \(R_3/F_3\) is abelian, while
\(\operatorname{Gal}(L_3F_3/F_3)=H_3\), since
\(L_3\cap F_3=K_3\). Its intersection with
\(R_3\) is fixed by \([H_3,H_3]=Z_3\).
Hence a nonzero element of \(Z_3\), paired with
the identity of \(R_3\), extends to the compositum.
Chebotarev gives infinitely many primes with this
Frobenius \cite[Theorem 8.31]{GhostMilneANT}.
Choose the least one above all the first-stage bounds.
It satisfies the original first-prime conditions,
so the unchanged A.44 induction gives an ordinary
pair with exactly the displayed original exponents.
All later choices are fixed by taking the least
eligible primes and finite residue choices.

Explicitly the added first-prime test is
\[
 \begin{gathered}
 q\equiv1\pmod{D_3},\qquad
 (-2)^{(q-1)/24}=1\pmod q,\qquad
 3^{(q-1)/3}=1\pmod q,\\
 \theta_q=(1-\zeta_qx_q)(1-\zeta_q^2x_q)^2,
 \quad x_q^3=3,\quad \operatorname{ord}(\zeta_q)=3,
 \qquad \theta_q^{(q-1)/3}\ne1.
 \end{gathered}                                      \tag{7}
\]
The same radical identity proves independence of the
choices of \(x_q,\zeta_q\). In particular
\(q_0\equiv1\pmod9\), and the point-three section
has Frobenius coordinates
\((0,0,c)\) modulo three, with \(c\ne0\).

Let \(s\) be any integral class-two section over this
pair. Normalize its abelian cocycles modulo three
to those of the point-three section. Their central
difference \(w\bmod3\) is a weight-two cocycle.
At each new support prime its inertia value is a
multiple of \(n_j\), by the integral formulas
\(h(T)=n_ju\) or \(-n_jv\); it is consequently
zero modulo three. At every zero-valuation prime,
integral torsion freeness forces inertia zero.
At two the difference satisfies
\((2^2-2)w(T)=0\), and this scalar is invertible
modulo three. Thus \(w\) factors through
\(G_{\mathbb Q,\{3,\infty\}}\).

Its first cohomology with \(\mathbb F_3(2)\)
coefficients is generated by the reduction of
\(\eta=(\chi^2-1)/24\). Indeed inflation embeds
integral first cohomology into the absolute group
\(\mathbb Z/3\) of
Proposition~\ref{prop:ghost-h1-weight2}; the unramified
cocycle \(\eta\) already generates that absolute
group, so inflation is an isomorphism here.
The integral second cohomology is
\(K_2(\mathbb Z)\{3\}=0\). The coefficient exact
sequence therefore identifies the finite first
cohomology with the quotient of integral first
cohomology by three. Consequently
\[
                    w=k\eta+\delta v\pmod3
\]
for some \(k\in\mathbb Z/3\) and an integral
zero-cochain \(v\) modulo three.
At \(q=q_0\), \(q\equiv1\pmod9\) gives
\[
 \eta(F_q)=(q^2-1)/24\equiv0\pmod3,
 \qquad (\delta v)(F_q)=(q^2-1)v\equiv0\pmod3.
\]
Thus every integral lift has
\(\alpha\equiv0\pmod3\) and
\(\gamma\equiv c\ne0\pmod3\) at this prime.
This explicitly accounts for the nonzero global
central freedom at the coefficient prime three.

Finally \(m=24\), \(n=0\), \(\beta=(q-1)u\),
and \(t=mu\). The same exact residue identity (6)
gives
\[
 Z_{D,q}=24(\gamma-q\alpha u)
                +(q^2-1)u\frac{24\cdot23}{2}.
\]
Its three-adic valuation is one, whereas
\(v_3(q^2-1)\ge2\). The class in
\(H^2(G_{\mathbb Q_q},\mathbb Z_3(3))
=\mathbb Z_3/(q^2-1)\mathbb Z_3\)
is nonzero, and remains so modulo nine. This excludes
every integral class-two lift at class three.
Class two exists for the ordinary pair by
Corollary~\ref{cor:ghost-class2-lift}, proving the
claimed first depth. As before, the finite detection
is applied to reductions of integral class-two
sections; no surjectivity onto arbitrary finite
coefficient section classes is asserted.
\end{proof}
